% =====================================================================
%  Bölüm 11 — Geçerlilik Teoremi
% =====================================================================
\chapter{Geçerlilik Teoremi}\label{bol:gecerlilik}

\begin{bolumhedefleri}
\begin{itemize}
  \item Bir ispat sistemini \Lean{}'de tümevarımsal bir yüklem olarak
        tanımlamayı,
  \item Nesne dilinde türetimler kurmayı ve onları ağaç olarak okumayı,
  \item Geçerlilik teoremini (``türetilebilen her şey doğrudur'') kural
        tümevarımıyla tam olarak ispatlamayı,
  \item Geçerlilikten tutarlılığı ve türetilemezlik sonuçlarını
        çıkarmayı,
  \item Yapısal kuralları (zayıflatma) ve bir ispat sisteminin
        anlambilimine göre neyi ``yakalayıp'' neyi yakalayamadığını.
\end{itemize}
\end{bolumhedefleri}

Bölüm~10'da mantığın iki yüzünden birini, anlambilimi tanımladık:
$\Gamma \models A$. Bu bölümde öteki yüzü, sözdizimini tanımlıyor ve
ikisi arasındaki ilk köprüyü kuruyoruz:
\[
  \Gamma \vdash A \quad\Longrightarrow\quad \Gamma \models A.
\]
Bu, kitabın ilk tam \emph{metamantık} teoremi olacak: Bir mantık sistemi
\emph{hakkında}, o sistemin dışından ispatlanan bir teorem.

% ---------------------------------------------------------------------
\section{Bir ispat sistemi, veri olarak}\label{sec:ispat-sistemi}

Bölüm~5'te doğal çıkarım kurallarını ardışık biçiminde, yani her
düğümde bağlamı açıkça yazarak ifade etmiştik. O kuralları şimdi
\Lean{}'de bir \emph{tümevarımsal yüklem} olarak yazıyoruz. Her kural bir
kurucu, her türetim ağacı ise bu yüklemin bir kanıtı olacak:

\leandosya[Klasik doğal çıkarım]{Bolum11}{turetilir}

\noindent \lc{Turetilir Γ A}, kâğıt üzerinde $\Gamma \vdash A$ diye
yazdığımız şeydir. Bağlamı (\lc{Γ}) bir liste olarak tutuyoruz;
yeni bir varsayım listenin başına eklenir. \Cref{fig:kural-sistemi},
on bir kuralın tamamını ardışık biçiminde gösteriyor.

\begin{figure}[htbp]
\centering
\setlength{\tabcolsep}{16pt}
\renewcommand{\arraystretch}{3.0}
\begin{tabular}{ccc}
\begin{prooftree}
  \hypo{A \in \Gamma}
  \infer1[Var]{\Gamma \vdash A}
\end{prooftree}
&
\begin{prooftree}
  \hypo{\Gamma \vdash A}
  \hypo{\Gamma \vdash B}
  \infer2[$\wedge$G]{\Gamma \vdash A \ove B}
\end{prooftree}
&
\begin{prooftree}
  \hypo{\Gamma \vdash A \ove B}
  \infer1[$\wedge$E$_{1,2}$]{\Gamma \vdash A \quad (B)}
\end{prooftree}
\\
\begin{prooftree}
  \hypo{\Gamma \vdash A \quad (B)}
  \infer1[$\vee$G$_{1,2}$]{\Gamma \vdash A \oveya B}
\end{prooftree}
&
\multicolumn{2}{c}{%
\begin{prooftree}
  \hypo{\Gamma \vdash A \oveya B}
  \hypo{A, \Gamma \vdash C}
  \hypo{B, \Gamma \vdash C}
  \infer3[$\vee$E]{\Gamma \vdash C}
\end{prooftree}}
\\
\begin{prooftree}
  \hypo{A, \Gamma \vdash B}
  \infer1[$\to$G]{\Gamma \vdash A \oise B}
\end{prooftree}
&
\begin{prooftree}
  \hypo{\Gamma \vdash A \oise B}
  \hypo{\Gamma \vdash A}
  \infer2[$\to$E]{\Gamma \vdash B}
\end{prooftree}
&
\begin{prooftree}
  \hypo{\Gamma \vdash \bot}
  \infer1[$\bot$E]{\Gamma \vdash A}
\end{prooftree}
\\
\multicolumn{3}{c}{%
\colorbox{teorembg}{\begin{prooftree}
  \hypo{{\sim}A, \Gamma \vdash \bot}
  \infer1[RAA]{\Gamma \vdash A}
\end{prooftree}\quad\raisebox{0.4em}{\textcolor{teoremframe}{\baslikfont\scriptsize\bfseries klasik kural}}}}
\\
\end{tabular}
\caption{\lc{Turetilir} yükleminin on bir kuralı. Her kural, \Lean{}
tanımındaki bir kurucuya karşılık gelir.}
\label{fig:kural-sistemi}
\end{figure}

\begin{onemlifikir}[Önemli Fikir: Türetim = kanıt terimi]
\lc{Turetilir Γ A} türünde bir \Lean{} terimi, bir türetim ağacının
ta kendisidir: Her kurucu bir düğüm, kurucunun argümanları o düğümün
çocuklarıdır. Bölüm~9'daki \lc{Cift} yükleminde \lc{Cift 4}'ün kanıtı
nasıl bir türetim ağacıysa, \lc{Turetilir [] peirce}'in kanıtı da bir
doğal çıkarım ağacıdır. Böylece ``türetim ağaçları üzerinde tümevarım'',
\Lean{}'de sıradan bir kural tümevarımına dönüşür.
\end{onemlifikir}

% ---------------------------------------------------------------------
\section{Nesne dilinde türetimler}\label{sec:nesne-turetim}

En basit türetim, özdeşlik yasasıdır: $\vdash A \oise A$.

\leandosya[Özdeşlik]{Bolum11}{ornek_ozdeslik}

\begin{center}
\begin{prooftree}
  \hypo{A \in [A]}
  \infer1[Var]{A \vdash A}
  \infer1[$\to$G]{\vdash A \oise A}
\end{prooftree}
\end{center}

\noindent Bu ispat ile Bölüm~2'deki \lc{theorem ozdeslik (P : Prop) : P → P}
arasındaki farka dikkat edin. Orada \Lean{}'in \emph{kendi} mantığında bir
önerme ispatlıyorduk (sığ gömme). Burada ise nesne dilinin bir formülünün,
nesne dilinin ispat sisteminde \emph{türetilebilir} olduğunu, üst dilde
ispatlıyoruz (derin gömme).

\subsection{Peirce yasasının türetimi}

Bölüm~4'te Peirce yasasının ancak klasik mantıkla ispatlanabildiğini
görmüştük. Nesne dilindeki ispat sistemimiz klasik olduğu için (RAA
kuralı sayesinde), Peirce yasası burada türetilebilir:

\leandosya[Peirce yasasının türetimi]{Bolum11}{ornek_peirce}

\begin{figure}[htbp]
\centering
\footnotesize
\begin{prooftree}
  \infer0[Var]{\Gamma_1 \vdash {\sim}p}
  \infer0[Var]{\Gamma_1 \vdash (p \oise q) \oise p}
  \infer0[Var]{\Gamma_2 \vdash {\sim}p}
  \infer0[Var]{\Gamma_2 \vdash p}
  \infer2[$\to$E]{\Gamma_2 \vdash \bot}
  \infer1[$\bot$E]{\Gamma_2 \vdash q}
  \infer1[$\to$G]{\Gamma_1 \vdash p \oise q}
  \infer2[$\to$E]{\Gamma_1 \vdash p}
  \infer2[$\to$E]{\Gamma_1 \vdash \bot}
  \infer1[\textcolor{teoremframe}{RAA}]{\Gamma_0 \vdash p}
  \infer1[$\to$G]{\vdash ((p \oise q) \oise p) \oise p}
\end{prooftree}

\vspace{6pt}
\normalsize
\begin{tabular}{r@{\;=\;}l}
$\Gamma_0$ & $[(p \oise q) \oise p]$ \\
$\Gamma_1$ & ${\sim}p :: \Gamma_0$ \\
$\Gamma_2$ & $p :: \Gamma_1$
\end{tabular}
\caption{Peirce yasasının türetim ağacı. Ağacın her düğümü, \Lean{}
ispatındaki bir \leantac{apply} ya da \leantac{exact} adımına karşılık
gelir. Klasik kural RAA yalnızca bir kez kullanılır.}
\label{fig:peirce-turetim}
\end{figure}

\noindent \Lean{} ispatındaki her \leantac{apply} adımı ağaçta bir
düğümü, her \lc{by decide} ise bir varsayımın bağlamda bulunduğunu
denetleyen bir hesabı temsil eder. \lc{by decide}'ın burada işe yaraması,
\lc{Formul} tipinin karar verilebilir eşitliğe sahip olmasındandır
(Bölüm~10'daki \leankw{deriving} \lc{DecidableEq}).

% ---------------------------------------------------------------------
\section{Geçerlilik teoremi}\label{sec:gecerlilik-teoremi}

\begin{teorem}{Geçerlilik}{gecerlilik}
Her $\Gamma$ ve $A$ için, $\Gamma \vdash A$ ise $\Gamma \models A$.
Yani türetilebilen her formül, varsayımlarının mantıksal sonucudur.
\end{teorem}

\subsection{İspatın fikri}

$\Gamma \vdash A$'nın bir türetimi verilsin. Türetim üzerinde kural
tümevarımı yapıyoruz. Her kural için şunu göstermeliyiz: Kuralın
\emph{öncülleri} geçerliyse (tümevarım hipotezleri), \emph{sonucu} da
geçerlidir. Başka bir deyişle, \textbf{her kural doğruluğu korur}
(\cref{fig:dogruluk-korunur}).

\begin{figure}[htbp]
\centering
\begin{tikzpicture}[
    dugum/.style={draw=bittiframe, fill=bittibg, thick, rounded corners=3pt,
                  font=\small, inner sep=4pt, minimum height=0.7cm},
    kural/.style={font=\baslikfont\scriptsize\bfseries, text=leantac},
    ok/.style={-{Stealth[length=2.4mm]}, thick, bittiframe},
  ]
  \node[dugum] (y1) at (-2.6,3.0) {$\Gamma \models A$};
  \node[dugum] (y2) at (-0.2,3.0) {$\Gamma \models A \oise B$};
  \node[dugum] (y3) at (2.6,3.0) {$\Gamma \models C$};
  \node[dugum] (o1) at (-1.4,1.5) {$\Gamma \models B$};
  \node[dugum] (k) at (0.6,0.0) {$\Gamma \models B \ove C$};
  \draw[ok] (y1) -- (o1); \draw[ok] (y2) -- (o1);
  \draw[ok] (o1) -- (k); \draw[ok] (y3) -- (k);
  \node[kural] at (-0.2,2.2) {$\to$E};
  \node[kural] at (1.4,0.85) {$\wedge$G};
  \node[font=\baslikfont\scriptsize, text=black!60, align=left, anchor=west] at (3.8,2.3)
       {yapraklar: Var\\(varsayım doğrudur)};
  \node[font=\baslikfont\scriptsize, text=black!60, align=left, anchor=west] at (3.8,0.6)
       {her kural:\\öncüller geçerliyse\\sonuç da geçerlidir};
  \draw[decorate, decoration={brace, amplitude=6pt, mirror}, thick, black!35]
       (-3.6,-0.4) -- (3.4,-0.4) node[midway, below=8pt, font=\baslikfont\small\bfseries, text=bittiframe]{doğruluk ağaç boyunca aşağı akar};
\end{tikzpicture}
\caption{Geçerlilik ispatının fikri. Yapraklardaki varsayımlar doğrudur;
her kural, öncüllerindeki doğruluğu sonucuna taşır. Dolayısıyla kökteki
formül de doğrudur.}
\label{fig:dogruluk-korunur}
\end{figure}

\noindent Her kuralın durumu, Bölüm~10'daki \lc{deger} tanımından
kolayca çıkar (\cref{tab:gecerlilik-durumlar}).

\begin{table}[htbp]
\centering
\small
\renewcommand{\arraystretch}{1.45}
\rowcolors{2}{tablosatir}{white}
\begin{tabularx}{\linewidth}{>{\raggedright\arraybackslash}p{1.6cm} X}
\rowcolor{tabloust}
\tablobaslik{Kural} & \tablobaslik{Neden doğruluğu korur? ($v$, $\Gamma$'yı doğrulayan bir değerleme)} \\
Var       & $A \in \Gamma$ ve $v$, $\Gamma$'nın her formülünü doğrular. \\
$\wedge$G & $A$ ve $B$ doğruysa \lc{true \&\& true = true}. \\
$\wedge$E & \lc{a \&\& b = true} ise \lc{a = true} ve \lc{b = true}. \\
$\vee$G   & $A$ doğruysa \lc{true || b = true}. \\
$\vee$E   & $A \oveya B$ doğruysa $A$ ya da $B$ doğrudur; o durumda genişletilmiş bağlam $v$ altında doğrudur ve tümevarım hipotezi $C$'yi verir. \\
$\to$G    & $A$ yanlışsa $A \oise B$ doğrudur; $A$ doğruysa $A{::}\Gamma$ doğrulanır ve hipotez $B$'yi verir. \\
$\to$E    & \lc{!a || b = true} ve \lc{a = true} ise \lc{b = true}. \\
$\bot$E   & $\bot$ hiçbir değerlemede doğru değildir: öncül \emph{boş yere} geçerlidir. \\
RAA       & $A$ yanlış olsaydı ${\sim}A$ doğru olurdu ve hipotez $\bot$'u doğrulardı. \textbf{Bu adım, $A$'nın değerinin yalnızca iki seçenekten biri olmasına dayanır.} \\
\end{tabularx}
\caption{Geçerlilik ispatının on bir durumu.}
\label{tab:gecerlilik-durumlar}
\end{table}

\subsection{Lean ispatı}

Önce, bağlama doğru bir formül eklemenin bağlamın doğruluğunu
bozmadığını gösteren küçük bir önsava ihtiyacımız var. Bu önsav,
bağlamı genişleten üç kuralda ($\vee$E, $\to$G, RAA) kullanılacak:

\leandosya[Bağlamı genişletmek]{Bolum11}{baglam_ekle}

\noindent Şimdi teoremin kendisi:

\leandosya[Geçerlilik teoremi]{Bolum11}{gecerlilik}

\noindent İspatta dikkat edilmesi gereken birkaç nokta:

\begin{itemize}
  \item \leantac{induction h with} satırındaki \lc{@varsayim Γ A hA}
        gibi örüntüler, kurucunun \emph{örtük} argümanlarına (\lc{Γ},
        \lc{A}, \lc{B}) ad vermemizi sağlar. Bu adlar, \leantac{show}
        ile hedefi açık yazabilmek için gereklidir.
  \item İkili bağlaçların durumlarında \leantac{simp}, Bölüm~10'daki
        \lc{Bool} işlemlerini (\lc{\&\&}, \lc{||}, \lc{!}) mantık
        bağlaçlarına çevirir: \lc{(a \&\& b) = true} hipotezi
        \lc{a = true ∧ b = true} olur.
  \item $\to$G ve RAA durumlarında \lc{deger v A}'nın \lc{true} mu
        \lc{false} mu olduğuna göre durum ayırırız. İki değerli
        anlambilimin gücü buradan gelir.
\end{itemize}

\begin{dikkat}[Dikkat: RAA neden geçerli?]
RAA kuralının geçerlilik ispatı, \lc{cases hA : deger v A} satırına
dayanır: Bir formülün değeri ya \lc{true} ya da \lc{false}'tur. Bu,
\Lean{}'in üst dilinde klasik mantık \emph{kullandığımız} anlamına
gelmez: \lc{Bool} üzerinde durum ayırmak tamamen yapıcıdır, çünkü
\lc{Bool}'un iki kurucusu vardır. Klasik olan, nesne dilinin \emph{anlambilimidir}: Her formüle
iki değerden birini veriyoruz. RAA bu anlambilime göre geçerlidir;
Bölüm~4'teki Kripke anlambilimine göre ise geçerli değildir.
\end{dikkat}

% ---------------------------------------------------------------------
\section{Sonuçlar}\label{sec:gecerlilik-sonuclar}

Geçerlilik teoreminin en önemli kullanımı, bir şeyin \emph{türetilemez}
olduğunu göstermektir. Bir formülü türetmeye çalışıp başarısız olmak hiçbir
şey kanıtlamaz; ama formülü yanlış yapan bir değerleme bulursak, geçerlilik
teoremi formülün hiçbir zaman türetilemeyeceğini garanti eder
(\cref{fig:karsi-model}).

\leandosya[Tutarlılık ve türetilemezlik]{Bolum11}{sonuclar}

\begin{figure}[htbp]
\centering
\begin{tikzpicture}[
    kutu/.style={draw=#1, fill=#1!8, thick, rounded corners=5pt,
                 text width=2.6cm, align=center, inner sep=6pt, font=\small},
    ok/.style={-{Stealth[length=3mm]}, very thick, #1},
  ]
  \node[kutu=hataframe] (k) at (0,0) {$A$'yı yanlış yapan\\bir değerleme $v$};
  \node[kutu=ikincil]  (s) at (5.2,0) {$\not\models A$};
  \node[kutu=anarenk]  (t) at (10.4,0) {$\nvdash A$\\{\scriptsize hiçbir türetim yok}};
  \draw[ok=hataframe] (k) -- node[above, font=\baslikfont\scriptsize]{tanım} (s);
  \draw[ok=teoremframe] (s) -- node[above, font=\baslikfont\scriptsize\bfseries, text=teoremframe]{geçerlilik} node[below, font=\baslikfont\scriptsize]{(kontrapozisyon)} (t);
\end{tikzpicture}
\caption{Türetilemezliği göstermenin yolu: bir karşı değerleme bulmak.
Geçerlilik teoreminin kontrapozisyonu, sonsuz sayıdaki olası türetimin
hiçbirinin $A$'yı vermeyeceğini garanti eder.}
\label{fig:karsi-model}
\end{figure}

\begin{teorem}{Tutarlılık}{tutarlilik-klasik}
Klasik doğal çıkarım tutarlıdır: $\nvdash \bot$.
\end{teorem}

\noindent Bölüm~5'te yapıcı mantığın tutarlılığını normalleştirme ve alt
formül özelliğiyle, yani \emph{sözdizimsel} yollarla göstermiştik. Burada
aynı sonuca, klasik sistem için, \emph{anlambilimsel} yoldan ulaşıyoruz:
$\bot$ hiçbir değerlemede doğru olmadığı için türetilemez. İki yöntem
birbirini tamamlar; ispat kuramında her ikisi de kullanılır.

% ---------------------------------------------------------------------
\section{Yapısal özellikler}\label{sec:yapisal}

Doğal çıkarımın kuralları belirli bir bağlaca bağlıdır. Bunların yanında
bir de bağlamın kendisiyle ilgili \emph{yapısal} özellikler vardır:

\begin{center}
\renewcommand{\arraystretch}{1.5}
\rowcolors{2}{tablosatir}{white}
\begin{tabular}{>{\raggedright\arraybackslash}p{2.8cm}
                >{\centering\arraybackslash}p{4.6cm}
                >{\raggedright\arraybackslash}p{4.6cm}}
\rowcolor{tabloust}
\tablobaslik{Özellik} & \tablobaslik{Kural} & \tablobaslik{Bizim sistemimizde} \\
zayıflatma  & $\Gamma \vdash A \;\Rightarrow\; B, \Gamma \vdash A$ & bir teorem (\lc{zayiflatma}) \\
değiştirme  & $A, B, \Gamma \vdash C \;\Rightarrow\; B, A, \Gamma \vdash C$ & zayıflatmanın özel hâli \\
daraltma    & $A, A, \Gamma \vdash C \;\Rightarrow\; A, \Gamma \vdash C$ & zayıflatmanın özel hâli \\
\end{tabular}
\end{center}

\noindent Bağlamı bir liste olarak tutsak da Var kuralı yalnızca
\emph{üyeliğe} baktığı için sıra ve tekrar önemli değildir. Bu yüzden
üç özelliğin üçü de tek bir teoremden çıkar: Bağlam genişletilirse
(daha doğrusu, bağlamın her elemanı yeni bağlamda da bulunuyorsa)
türetilebilirlik korunur.

\leandosya[Zayıflatma]{Bolum11}{zayiflatma}

\noindent \leantac{induction h generalizing Δ} satırına dikkat edin. Bağlamı
genişleten kurallarda (\lc{veyaE}, \lc{iseG}, \lc{raa}) tümevarım
hipotezini \emph{farklı} bir \lc{Δ} ile, yani \lc{A :: Δ} ile kullanmamız
gerekir. Bunun için tümevarımdan önce \lc{Δ}'yı genelleştiririz; böylece
tümevarım hipotezi her \lc{Δ} için geçerli olur. Bölüm~7'deki ``önermeyi
güçlendirme'' tekniğinin bir başka örneği budur.

% ---------------------------------------------------------------------
\section{Geçerlilik neyi söylemez?}\label{sec:gecerlilik-sinir}

Geçerlilik teoremi bir yöndedir: Türetilebilen her şey doğrudur. Ters
yön, yani doğru olan her şeyin türetilebilir olması, \emph{tamlık}
teoremidir ve ispatı çok daha zordur (\cref{bol:tamlik}).

Bu noktada RAA kuralının rolü önemlidir. RAA'yı sistemden çıkarırsak,
geri kalan on kural \emph{yapıcı} doğal çıkarımı oluşturur. Bu sistem de
iki değerli anlambilime göre geçerlidir: Ne de olsa RAA'sız her türetim,
RAA'lı sistemde de bir türetimdir. Ama tam değildir: Peirce yasası bir
totolojidir, yine de RAA olmadan türetilemez. Bunu ispatlamak için iki
değerli anlambilim yetmez; Kripke modellerine ihtiyaç vardır
(\ciltiki{2}).

\begin{onemlifikir}[Önemli Fikir: İspat sistemi ile anlambilimin uyumu]
Bir ispat sistemi ile bir anlambilim arasında iki ilişki olabilir:
Sistem anlambilime göre \emph{geçerli} olabilir (çok fazla şey
türetmez) ve \emph{tam} olabilir (yeterince şey türetir). Klasik doğal
çıkarım, iki değerli anlambilime göre hem geçerli hem tamdır. Yapıcı
doğal çıkarım ise iki değerli anlambilime göre geçerli ama tam
değildir; Kripke anlambilimine göre hem geçerli hem tamdır. Her mantığın
kendine uygun anlambilimi vardır.
\end{onemlifikir}

% ---------------------------------------------------------------------
\begin{arastirmanotu}
\textbf{Yapısal kuralların ötesi.} Zayıflatma ve daraltma kurallarının
\emph{kısıtlandığı} mantıklara \emph{alt yapısal mantıklar} denir.
Jean-Yves Girard'ın \emph{doğrusal mantığı} \citep{girard1987}, her
varsayımın tam olarak bir kez kullanılmasını ister ve kaynak kullanımını
modellemek için bilgisayar biliminde önemli bir rol oynar; programlama
dillerindeki kaynak ve sahiplik tiplerinin kuramsal arka planında yer alır. Bu tür
mantıkları \Lean{}'de formalize etmek için bağlamı liste yerine çoklu
küme (\emph{multiset}) olarak tutmak gerekir.

\textbf{Metamantığın formalizasyonu.} Önermeler mantığı için geçerlilik
ve tamlık teoremlerinin formalizasyonu, ispat asistanları için klasik
bir alıştırmadır; birinci mertebe mantık ve Gödel teoremleri ise çok daha
zordur. John Harrison'ın HOL Light'taki çalışmaları \citep{harrison2009}
ve Lean~4'teki Foundation projesi \citep{foundation, saitou2026} bu
alanın önemli örnekleridir.
\end{arastirmanotu}

\begin{bolumozeti}
\begin{itemize}
  \item Klasik doğal çıkarım, \Lean{}'de on bir kurucusu olan
        tümevarımsal bir yüklemdir: \lc{Turetilir Γ A} ($\Gamma \vdash A$).
        Bir türetim, bu yüklemin bir kanıtıdır.
  \item Nesne dilinde türetimler \leantac{apply} ile kurucuları
        uygulayarak kurulur; bağlam üyelikleri \lc{by decide} ile
        denetlenir.
  \item Geçerlilik teoremi ($\Gamma \vdash A \Rightarrow \Gamma \models A$)
        kural tümevarımıyla ispatlanır: Her kural doğruluğu korur.
  \item Geçerlilik, türetilemezliği karşı değerlemelerle göstermeyi ve
        sistemin tutarlılığını ($\nvdash \bot$) sağlar.
  \item Zayıflatma, değiştirme ve daraltma tek bir teoremden çıkar;
        ispatında bağlam genelleştirilir.
  \item RAA, iki değerli anlambilime göre geçerlidir; RAA'sız yapıcı
        sistem geçerlidir ama bu anlambilime göre tam değildir.
\end{itemize}
\end{bolumozeti}

\begin{alistirmalar}
\alistirma{1} $p \vdash p$. (\lc{a11\_1})

\alistirma{1} $K$ aksiyomu: her $A, B$ için $\vdash A \oise B \oise A$.
Bağlamın hangi elemanını kullandığınıza dikkat edin. (\lc{a11\_2})

\alistirma{2} $p \ove q \vdash q \ove p$. (\lc{a11\_3})

\alistirma{2} Hipotetik kıyas, nesne dilinde:
$p \oise q,\; q \oise r \vdash p \oise r$. Türetim ağacını da çizin ve
Bölüm~2'deki \lc{hipotetik\_kiyas} ile karşılaştırın. (\lc{a11\_4})

\alistirma{2} $\nvdash p \oise q$. Hangi değerlemeyi seçtiniz?
(\lc{a11\_5})

\alistirma{2} $p \oveya q \nvdash p$. (\lc{a11\_6})

\alistirma{3} Üçüncü hâlin olanaksızlığı: $\vdash p \oveya {\sim}p$.
Bölüm~3'teki \lc{em\_yanlis\_olamaz} ispatıyla karşılaştırın: RAA hangi
adımın yerini tutuyor? (\lc{a11\_7})

\alistirma{3} Çift değilleme eleme: ${\sim}{\sim}A \vdash A$. (\lc{a11\_8})

\alistirma{2} (Kâğıt-kalem) \Cref{tab:gecerlilik-durumlar}'deki $\vee$E
satırını ayrıntılı bir matematiksel ispata dönüştürün.

\alistirma{3} (Proje) RAA kuralı olmayan bir \lc{TuretilirY} yüklemi
tanımlayın. \lc{TuretilirY Γ A → Turetilir Γ A} önermesini ispatlayın ve
buradan yapıcı sistemin geçerliliğini tek satırda çıkarın.

\alistirma{3} (Tartışma) Bir arkadaşınız, geçerlilik teoreminin
\lc{simp} kullanmadan ispatlanmasını istiyor. Hangi durumlar zorlaşır?
\lc{Bool} yerine \lc{Prop} değerli bir anlambilim (\lc{deger v A : Prop})
kullansaydık RAA durumunda ne değişirdi?
\end{alistirmalar}
